%% ma: L12-B08-0048
%% lop: 12
%% chuong: 2
%% bai: 8
%% dang: 12.8.40
%% loai: TN4
%% muc: VD
%% dap_an: "B"
%% nguon: "LaTeX giáo viên gửi 21/10/2026 (M05, Vectơ và hệ trục tọa độ trong không gian, 107 câu), Phần I Câu 61"
%% kien_thuc: "Bài 7 (tọa độ điểm), Bài 8 (trọng tâm tam giác)"
%% trang_thai: chinh-thuc
%% ghi_chu: "Giáo viên đã duyệt 21/10/2026. Tọa độ các khối lấy theo hình (gốc O ở đỉnh khuất, các trục dọc theo ba cạnh). Hình dựng lại theo hình gốc"
\begin{ex}
Trong không gian với hệ tọa độ $Oxyz$, một hình lập phương có độ dài cạnh bằng $4$ được chia thành $8$ hình lập phương nhỏ có độ dài cạnh bằng $2$. Xét ba hình lập phương nhỏ ghi chữ $A$, $B$, $C$ như hình vẽ, bên trong mỗi hình này có một mặt cầu bán kính bằng $1$ nội tiếp. Gọi $(p;q;r)$ là tọa độ trọng tâm của tam giác tạo bởi tâm của $3$ mặt cầu đó. Giá trị của $p+q+r$ bằng
\begin{center}
\begin{tikzpicture}[x={(-.6cm,-.45cm)},y={(1cm,0cm)},z={(0cm,1cm)},scale=.65]
  \draw (4,0,0)--(4,4,0)--(4,4,4)--(4,0,4)--cycle (4,2,0)--(4,2,4) (4,0,2)--(4,4,2);
  \draw (4,4,0)--(0,4,0)--(0,4,4)--(4,4,4) (2,4,0)--(2,4,4) (4,4,2)--(0,4,2);
  \draw (4,0,4)--(0,0,4)--(0,4,4) (2,0,4)--(2,4,4) (4,2,4)--(0,2,4);
  \draw[truc] (4,0,0)--(6.5,0,0) node[below]{$x$};
  \draw[truc] (0,4,0)--(0,6,0) node[below]{$y$};
  \draw[truc] (0,0,4)--(0,0,6) node[left]{$z$};
  \node at (4,1,3) {$A$}; \node at (4,3,1) {$B$}; \node at (1,4,1) {$C$};
\end{tikzpicture}
\end{center}
\choice{$6$}{\True $\dfrac{19}3$}{$\dfrac{20}3$}{$7$}
\loigiai{Hình lập phương lớn chiếm $[0;4]^3$. Tâm mặt cầu nội tiếp mỗi khối nhỏ là tâm khối đó:
khối $A$ ($x,z\in[2;4]$, $y\in[0;2]$) có tâm $(3;1;3)$; khối $B$ ($x,y\in[2;4]$, $z\in[0;2]$) có tâm $(3;3;1)$; khối $C$ ($x,z\in[0;2]$, $y\in[2;4]$) có tâm $(1;3;1)$.

Trọng tâm $\left(\dfrac73;\dfrac73;\dfrac53\right)$, nên $p+q+r=\dfrac{19}3$.}
\end{ex}
